\subsection{仿射 Springer 纤维}

\subsubsection{仿射 Grassmann 簇中的仿射 Springer 纤维} 对任意 $k$ 代数 $R$，把 $\frg\otimes_{k}R$ 记为 $\frg(R)$。特别地，$\frg(F)=\frg\otimes_{k}F$ 是圈群 $LG$ 的 Lie 代数。对 $g\in LG$，用 $\Ad(g)$ 记它在 $\frg(F)$ 上的伴随作用。

设 $\gamma\in\frg(F):=\frg\otimes_{k}F$ 是正则半单元\footnote{这里涉及的 Lie 代数是域 $F$（它不是代数闭域）上的。称 $\gamma\in\frg(F)$ 为正则半单元，指它作为 $\frg(\overline{F})$ 中的元素正则半单，参见 \S\ref{sss:Gr alt} 的脚注。等价地说，$\gamma$ 正则半单当且仅当它在 $\frc(F)$ 中的像落在 $\frc^{\rs}(F)$ 中。}。考虑 $\Gr_{G}$ 的子函子：它在 $k$ 代数 $R$ 上的取值由
\begin{equation}\label{aff Spr}
\tX_{\gamma}(R)=\{[g]\in \Gr_{G}(R)|\Ad(g^{-1})\gamma\in\frg(R\tl{t})\}
\end{equation}
给出。则 $\tX_{\gamma}$ 是 $\Gr_{G}$ 的闭子归纳概形。令 $\sX_{\gamma}=\tX^{\red}_{\gamma}$ 为 $\tX_{\gamma}$ 底层的约化归纳概形。称 $\sX_{\gamma}$ 为 $\gamma$ 在仿射 Grassmann 簇 $\Gr_{G}$ 中的{\em 仿射 Springer 纤维}。

\subsubsection{用格给出的另一套定义}\label{sss:ASF Gr GLn} 考虑 $G=\GL_{n}$ 的情形。设 $\gamma\in\frg(F)=\gl_{n}(F)$ 是正则半单矩阵。按 \S\ref{sss:Gr} 的做法，把 $\Gr_{G}$ 等同于 $F^{n}$ 中格的模空间；更精确地说，$\Gr_{G}(R)$ 是 $R\lr{t}^{n}$ 中格的集合。此时 $\tX_{\gamma}(R)$ 可以等同于那些满足 $\gamma \L\subset \L$ 的格 $\L\subset R\lr{t}^{n}$，即那些在 $R\lr{t}^{n}$ 的自同态 $\gamma$ 下稳定的格。

当 $G=\SL_{n}$ 时，$\Gr_{G}(R)$ 分类 $R\lr{t}^{n}$ 中满足 $[\L:R\tl{t}^{n}]=0$ 的格 $\L$。此时的仿射 Springer 纤维 $\tX_{\gamma}$ 仍由同样的条件 $\gamma \L\subset \L$ 刻画。

当 $G=\Sp_{2n}$ 时，$\Gr_{G}(R)$ 分类 $R\lr{t}^{2n}$ 中满足 $\L^{\vee}=\L$ 的格 $\L$，参见 \S\ref{sss:Sp parahoric}。此时的仿射 Springer 纤维 $\tX_{\gamma}$ 也由同样的条件 $\gamma \L\subset \L$ 刻画。

下面给出仿射 Springer 纤维最简单的几个例子。

\subsubsection{}\label{sss:reduction rs} 设 $\gamma\in\frt(\calO_{F})$，且其约化 $\ov\gamma\in\frt$ 正则半单。对每个余权 $\l:\Gm\to T$，元素 $t^{\l}:=\l(t)\in T(F)$ 给出点 $[t^{\l}]\in \Gr_{G}(k)$；由于 $\Ad(t^{-\l})\gamma=\gamma\in \frg(\calO_{F})$，该点落在 $\tX_{\gamma}$ 中。约化归纳概形 $\sX_{\gamma}$ 实际上就是离散集 $\{[t^{\l}]\}$，它与 $\xcoch(T)$ 一一对应。更典范地说，圈群 $LT$ 通过它在 $\Gr_{G}$ 上的左平移作用而作用在 $\tX_{\gamma}$ 上。此作用经由商 $\Gr_{T}=LT/L^{+}T$ 分解，并把 $\tX_{\gamma}$ 实现为一个 $\Gr_{T}$-主齐性空间。

\subsubsection{}\label{P1 chain} 考虑 $G=\SL_{2}$ 与 $\gamma=\mat{t}{0}{0}{-t}$ 的情形。此时 $\sX_{\gamma}$ 是 $\PP^{1}$ 的无穷长链。更精确地说，对每个 $n\in\ZZ$，考虑 $\Gr_{G}$ 的子概形 $C_{n}$，它分类满足 $t^{n}\calO_{F}\oplus t^{-n+1}\calO_{F} \subset \L\subset t^{n-1}\calO_{F}\oplus t^{-n}\calO_{F}$ 的格 $\L\subset F^{2}$。则 $C_{n}\cong\PP^{1}$。于是 $\sX_{\gamma}=\cup_{n\in\ZZ}C_{n}$ 是 $\PP^{1}$ 的无穷长链：分支 $C_{n}$ 与 $C_{n+1}$ 交于一点 $t^{-n}\calO_{F}\oplus t^{n}\calO_{F}$，此外诸分支互不相交。

下面给出一种计算 $\sX_{\gamma}$ 的 $k$ 点的方法。使用 $G(F)$ 的{\em Iwasawa 分解}：
\begin{equation*}
G(F)=\bigsqcup_{n\in\ZZ}N(F)\mat{t^{n}}{0}{0}{t^{-n}}G(\calO_{F}).
\end{equation*}
按此分解，$\Gr_{G}$ 中任一点都可由
\begin{equation}\label{Iwasawa elt}
g=\mat{1}{x}{0}{1}\mat{t^{n}}{0}{0}{t^{-n}}
\end{equation}
表示，其中 $x\in F$、$n\in\ZZ$ 唯一，且 $x$ 在 $F/t^{2n}\calO_{F}$ 中有确定的像。由于
\begin{equation*}
\Ad(g^{-1})\gamma=g^{-1}\gamma g=\mat{t}{2t^{1-2n}x}{0}{-t}
\end{equation*}
条件 $\Ad(g^{-1})\gamma\in\frg(\calO_{F})$ 等价于要求 $x\in t^{2n-1}\calO_{F}$。因此 $\sX_{\gamma}(k)=\sqcup_{n\in\ZZ}Y_{n}$，其中 $Y_{n}$ 由形如 \eqref{Iwasawa elt} 且 $x\in t^{2n-1}\calO_{F}/t^{2n}\calO_{F}$ 的元素组成。于是每个 $Y_{n}$ 都可等同于 $k$。容易验证 $Y_{n}\subset C_{n}(k)$。

\subsubsection{}\label{SL2 ram dim 1} 考虑 $G=\SL_{2}$ 与 $\gamma=\mat{0}{t^{2}}{t}{0}$ 的情形。此时 $\sX_{\gamma}$ 恰由那些满足 $t\calO_{F}\oplus \calO_{F}\subset\L\subset\calO_{F}\oplus t^{-1}\calO_{F}$ 的格 $\L\in\Gr_{G}$ 组成。因此 $\sX_{\gamma}\cong\PP^{1}$。这些计算的细节留给读者，见习题 \ref{exII:SL2}。

\subsubsection{共轭不变性} 若 $\gamma,\gamma'\in\frg(F)$ 满足 $\Ad(g)\gamma=\gamma'$（$g\in G(F)$），则 $\Gr_{G}$ 上的左乘 $g$ 限制为同构 $\tX_{\gamma}\cong\tX_{\gamma'}$，从而也有 $\sX_{\gamma}\cong\sX_{\gamma'}$。因此 $\sX_{\gamma}$ 的同构类型在 $\gamma$ 的 $G(F)$-共轭下不变。回忆有映射 $\chi:\frg\to \frc:=\frg\sslash G\cong\frt\sslash W$。对正则半单点 $a\in\frc^{\rs}(F)$，纤维 $\chi^{-1}(a)$ 是单个 $G(F)$-共轭类（这里用到 $F$ 的剩余域代数闭）。因此 $\sX_{\gamma}$ 的同构类型只依赖于 $a=\chi(\gamma)\in\frc^{\rs}(F)$。

与 Springer 纤维不同，$\sX_{\gamma}$ 对某些 $\gamma$ 可以是空集。仿射 Springer 纤维 $\sX_{\gamma}$ 非空，当且仅当 $a=\chi(\gamma)\in\frc(\calO_{F})$。事实上，若 $gG(\calO_{F})\in\sX_{\gamma}(k)$，则 $\Ad(g^{-1})\gamma\in\frg(\calO_{F})$，故 $\chi(\gamma)=\chi(\Ad(g^{-1})\gamma)\in\frc(\calO_{F})$。反之，$\chi$ 有 Kostant 截面 $\ep: \frc\to \frg$，它把 $\frc$ 等同于 $e+\frg_{f}$，其中 $(e,h,f)$ 是 $\frg$ 中的正则 $\sl_{2}$ 三元组。因此对任意 $a\in\frc(\calO_{F})\cap\frc^{\rs}(F)$，有 $\ep(a)\in\frg(\calO_{F})$，于是 $\sX_{\ep(a)}$ 包含 $\Gr_{G}$ 中的单位陪集，从而非空；再由 $\sX_{\gamma}$ 与 $\sX_{\ep(\chi(\gamma))}$ 同构，$\sX_{\gamma}$ 也非空。

对 $a\in\frc(\calO_{F})\cap\frc^{\rs}(F)$，也把 $\sX_{\ep(a)}$ 记为 $\sX_{a}$。上面的讨论表明，所有 Springer 纤维 $\sX_{\gamma}$ 都同构于 $\sX_{a}$，其中 $a=\chi(\gamma)$。

\subsubsection{parahoric 版本} 对每个 parahoric 子群 $\bP\subset LG$，可以用 \eqref{aff Spr} 的条件的类似物（把 $\frg(R\tl{t})$ 换成 $(\Lie\bP)\widehat{\otimes}_{k}R$）定义闭子归纳概形 $\tX_{\bP,\gamma}\subset\Fl_{\bP}$。约化归纳概形 $\sX_{\bP,\gamma}=\tX_{\bP,\gamma}^{\red}\subset\Fl_{\bP}$ 称为 $\gamma$ 的{\em $\bP$ 型仿射 Springer 纤维}。特别地，当 $\bI$ 是 $LG$ 的 Iwahori 子群时，把 $\sX_{\bI,\gamma}$ 记为 $\sY_{\gamma}$。

设 $\bP\subset\bQ$ 是 $LG$ 的两个 parahoric 子群。自然投影 $\Fl_{\bP}\to \Fl_{\bQ}$ 诱导映射 $\sX_{\bP,\gamma}\to \sX_{\bQ,\gamma}$。特别地，有映射 $\sY_{\gamma}\to \sX_{\gamma}$。

\subsection{仿射 Springer 纤维上的对称性}
对 Springer 纤维 $\sB_{e}$，中心化子 $C_{G}(e)$ 作用在它上面。本节研究仿射 Springer 纤维上的类似结构。

\subsubsection{中心化子的作用}\label{sss:centralizer} 设 $G_{\gamma}$ 是 $\gamma$ 在 $G_{F}$（$G$ 经基变换得到的 $F$ 上代数群）中的中心化子。则 $G_{\gamma}$ 是 $F$ 上代数群。由于 $\gamma$ 正则半单，$G_{\gamma}$ 是 $F$ 上的环面。可以定义 $G_{\gamma}$ 的圈群 $LG_{\gamma}$ 为 $k$ 代数上的函子 $R\mapsto G_{\gamma}(R\lr{t})$。

我们断言 $LG_{\gamma}$ 作用在归纳概形 $\tX_{\gamma}$ 上。这可以在 $R$ 点上看出。设 $h\in LG_{\gamma}(R)=G_{\gamma}(R\lr{t})$ 与 $[g]\in \tX_{\gamma}(R)$。则陪集 $[hg]\in \Gr_{G}(R)$ 仍满足
\begin{equation*}
\Ad((hg)^{-1})\gamma=\Ad(g^{-1})\Ad(h^{-1})\gamma=\Ad(g^{-1})\gamma\in\frg(R\tl{t}),
\end{equation*}
这里用到 $h$ 在 $\gamma$ 的中心化子中。因此 $[hg]\in \tX_{\gamma}(R)$。映射 $[g]\mapsto [hg]$（$h\in LG_{\gamma}$，$[g]\in\tX_{\gamma}$）定义了 $LG_{\gamma}$ 在 $\tX_{\gamma}$ 上的作用。它诱导出 $LG_{\gamma}$ 在约化结构 $\sX_{\gamma}$ 上的作用。

\subsubsection{分裂情形} 考虑 $\gamma\in\frt^{\rs}(F)$ 的情形。此时 $G_{\gamma}=T\otimes_{k}F$，且
\begin{equation*}
LG_{\gamma}=LT\cong\xcoch(T)\otimes_{\ZZ}L\Gm,
\end{equation*}
其中 $L\Gm$ 是乘法群 $\Gm$ 的圈群。对任意 $k$ 代数 $R$，$L\Gm(R)=R\lr{t}^{\times}$。容易看出，元素 $a=\sum_{i}a_{i}t^{i}\in R\lr{t}$ 可逆，当且仅当 $a$ 以有限多个幂零系数开头、且第一个非幂零系数在 $R$ 中可逆。当 $R$ 约化时，$a$ 的首项系数必在 $R$ 中可逆，于是 $R\lr{t}^{\times}=t^{\ZZ}\cdot R\tl{t}^{\times}$，而 $R\tl{t}^{\times}=\{(a_{0},a_{1},\cdots)|a_{0}\in R^{\times}, a_{i}\in R, \forall i\geq 1\}$。由此看出约化归纳概形 $(L\Gm)^{\red}\cong \ZZ\times L^{+}\Gm$，且作为概形 $L^{+}\Gm\cong \Gm\times\AA^{\NN}$，其中 $\AA^{\NN}=\Spec k[x_{1},x_{2},\cdots]$。因此当 $\gamma\in\frt^{\rs}(F)$ 时，$(LT)^{\red}\cong\xcoch(T)\times L^{+}T$，而 $L^{+}T$ 是无穷型的仿射概形。

\subsubsection{格 $\L_{\gamma}$}\label{sss:L gamma} 对一般的正则半单 $\gamma\in\frg(F)$，记 $\xcoch(G_{\gamma}):=\Hom_{F}(\Gm,G_{\gamma})$ 为环面 $G_{\gamma}$ 的 $F$-有理余权格。每个 $\l\in\xcoch(G_{\gamma})$ 可视为定义在 $F$ 上的同态 $\Gm\to G_{\gamma}$，于是可以谈论元素 $\l(t)$。对应 $\l\mapsto \l(t)$ 定义了单同态
\begin{equation*}
\xcoch(G_{\gamma})\incl G_{\gamma}(F),
\end{equation*}
其像记为 $\L_{\gamma}$。可以证明商 $\L_{\gamma}\bs (LG_{\gamma})^{\red}$ 是仿射概形，并且是若干 $\Gm^{a}\times \AA^{\NN}$（$a$ 为某整数）的有限不交并。

\subsubsection{$G=\GL_{n}$ 的情形}\label{sss:GLn centralizer} 沿用 \S\ref{sss:ASF Gr GLn} 的设定。假设 $\chark>n$。则 $\gamma$ 的特征多项式 $P(x)=x^{n}+a_{1}x^{n-1}+\cdots+a_{n}\in F[x]$ 可分。于是 $F$ 代数 $F[x]/(P(x))$ 是域之积 $F_{1}\times \cdots \times F_{m}$，且 $\sum_{i=1}^{m}[F_{i}:F]=n$。每个域扩张 $F_{i}/F$ 由添加 $P(x)$ 的不可约因子 $P_{i}(x)$ 的一个根得到，而由 $\chark>n$ 知 $F_{i}$ 必形如 $k\lr{t^{1/e_{i}}}$。此时中心化子 $G_{\gamma}$ 同构于{\em 诱导环面}之积：
\begin{equation*}
G_{\gamma}\cong \prod_{i=1}^{m}\Res^{F_{i}}_{F}\Gm.
\end{equation*}
有 $\xcoch(G_{\gamma})\cong \ZZ^{m}$，而映射 $\xcoch(G_{\gamma})\to G_{\gamma}(F)$ 由
\begin{equation*}
\ZZ^{m}\ni(d_{1},\cdots, d_{m})\mapsto (t^{d_{1}},\cdots, t^{d_{m}})\in F^{\times}_{1}\times\cdots\times F^{\times}_{m}
\end{equation*}
给出。商 $\L_{\gamma}\bs G_{\gamma}(F)$ 同构于 $\prod_{i=1}^{m}F^{\times}_{i}/t^{\ZZ}$。由于每个 $F_{i}$ 都同构于 $k\lr{t^{1/e_{i}}}$，有正合列 $1\to \calO^{\times}_{F_{i}}\to F^{\times}_{i}/t^{\ZZ}\to \ZZ/e_{i}\ZZ\to 0$，从而商 $\L_{\gamma}\bs (LG_{\gamma})^{\red}$ 以有限指数包含群概形 $\prod_{i=1}^{m}L^{+}_{F_{i}}\Gm$。这里 $L^{+}_{F_{i}}\Gm$ 作为概形同构于 $L^{+}\Gm$，只是把素元改名为 $t^{1/e_{i}}$。

另一种做法是：对每个 $i$ 固定素元 $t_{i}\in F_{i}$（例如取 $t_{i}=t^{1/e_{i}}$），并令 $\wt\L_{\gamma}=t_{1}^{\ZZ}\times\cdots\times t^{\ZZ}_{m}\subset \prod_{i}F_{i}^{\times}=G_{\gamma}(F)$。格 $\wt\L_{\gamma}$ 在计算轨道积分时有用，参见 \S\ref{sss:GLn orb coho}。

\subsubsection{$G=\SL_{2}$ 的情形} 设 $G=\SL_{2}$ 与 $\gamma=\mat{0}{t^{n}}{1}{0}$，其中 $n\geq 1$ 为奇数。则 $G_{\gamma}(F)$ 由满足 $a^{2}-t^{n}b^{2}=1$ 的矩阵 $\mat{a}{bt^{n}}{b}{a}$（$a,b\in F$）组成。注意此方程迫使 $a,b\in\calO_{F}$，故 $G_{\gamma}(F)=G_{\gamma}(\calO_{F})$。环面 $G_{\gamma}$ 非分裂，在二次扩张 $E=F(t^{1/2})$ 上分裂。此时格 $\L_{\gamma}=\Hom_{F}(\Gm,G_{\gamma})=0$。把 $a$ 与 $b$ 分别写成 $a=\sum_{i\geq0}a_{i}t^{i}$ 与 $b_{i}=\sum_{i\geq 0}b_{i}t^{i}$\zy{原文此处写作「$b_{i}=\sum_{i\geq0}b_{i}t^{i}$」，下标 $i$ 显系误置，应为「$b=\sum_{i\geq0}b_{i}t^{i}$」。译文照录，并登记为勘误 E2。［来源：原文 \texttt{ZhiweiYunPCMI.tex} 第 1431 行；下文「一旦 $b$ 与 $a_{0}$ 固定」亦要求以 $b$ 为对象。］}，可以看到 $a_{0}=\pm1$，并且一旦 $b$ 与 $a_{0}$ 固定，$a$ 的更高次系数就可用 $(1+t^{n}b^{2})^{1/2}$ 的 Taylor 展开唯一地解出。因此 $(LG_{\gamma})^{\red}\cong L^{+}G_{\gamma}$ 同构于 $\{\pm1\}\times \AA^{\NN}$，同构由 $(a,b)\mapsto (a_{0},b_{0},b_{1},\cdots)$ 给出。

\subsubsection{仿射 Springer 纤维上的对称性}\label{sss:sym ASF} 关于 $LG_{\gamma}$ 在 $\tX_{\gamma}$ 上的作用，Ng\^o 找到了更精确的结论，即该作用经由一个典范的有限维商分解。下面按 \cite[\S3.3]{NgoFL} 概述此事。设 $a=\chi(\gamma)\in\frc(F)^{\rs}$ 是 $\gamma$ 在 $\chi:\frg\to \frc$ 下的像。假定 $a\in\frc(\calO_{F})$，否则 $\sX_{\gamma}$ 为空。

存在 $\frc$ 上的光滑仿射群概形 $J$，称为{\em 正则中心化子群概形}。它的刻画性质是：它经由 $\chi$ 拉回到 $\frg$ 所得的 $\chi^{*}J$ 映入 $\frg$ 上的万有中心化子群概形 $I$，且该映射在正则轨迹 $\frg^{\reg}$ 上是同构。令 $J_{a}$ 为 $J$ 在映射 $a:\Spec\calO_{F}\to \frc$ 下的拉回。则 $J_{a}$ 是 $\calO_{F}$ 上的光滑仿射群概形，其 $F$-纤维是环面 $G_{\gamma}$（即 $J_{a}$ 是 $G_{\gamma}$ 在 $\calO_{F}$ 上的整模型）。可以构造 $J_{a}$ 的正圈群 $L^{+}J_{a}$，以及它的仿射 Grassmann 簇 $P_{a}:=LG_{\gamma}/L^{+}J_{a}$（也称为{\em 局部 Picard 群}）。约化群概形 $P^{\red}_{a}$ 是有限维的，且局部有限型。Ng\^o 证明：$LG_{\gamma}$ 在 $\tX_{\gamma}$（从而在 $\sX_{\gamma}$）上的作用经由局部 Picard 群 $P_{a}$ 分解，且不能经由更小的商分解。相关陈述见 \S\ref{dim formula discussion}。
